The position of the Sun on June 19, 2008:
$\delta \approx +23^\circ\!.5$, $\alpha \approx 6^{\mathrm{h}} = 90^\circ$.

At full moon phase, the declination of the Moon is between
$\delta \approx -18^\circ\!.5$ and $\delta \approx -28^\circ\!.5$ depending on
the position of node. \hfill(20\%)

% FIGURE NEEDED: celestial sphere with zenith Z, south celestial pole P, Moon M, horizon points N, E, S, W and the celestial equator (2008 TH solutions, p. 2)

Spherical triangle $ZPM$, $Z$ = zenith, $P$ = south celestial pole and
$M$ = moon.
\begin{align*}
  ZP &= m = (90^\circ - 6^\circ 49') = 83^\circ 11';\quad
  ZM = p = 90^\circ\ \text{and} \\
  PM &= z = 90^\circ - 23^\circ 30' = 66^\circ 30'
\end{align*}
The greatest possible value of $PM = z = 90^\circ - 18^\circ 30' = 71^\circ 30'$
and the smallest possible value of $PM = z = 90^\circ - 28^\circ 30'
= 61^\circ 30'$.
\[
  \cos p = \cos m \cos z + \sin m \sin z \cos\angle P
\]
\hfill(20\%)

The maximum value of $PM$
\begin{align*}
  \cos 90^\circ &= \cos 83^\circ 11' \cos 71^\circ 30'
    + \sin 83^\circ 11' \sin 71^\circ 30' \cos\angle P \\
  \cos\angle P &= -\cot 83^\circ 11' \cot 71^\circ 30'
    = -0.1195 \times 0.3346 \\
  &= -0.0401 \\
  \angle P &= 92^\circ 17' 53'' = 6^{\mathrm{h}} 9^{\mathrm{m}} 11^{\mathrm{s}} \\
  2\angle P &= 12^{\mathrm{h}} 18^{\mathrm{m}} 22^{\mathrm{s}}
\end{align*}
\hfill(30\%)

The minimum value of $PM$
\begin{align*}
  \cos 90^\circ &= \cos 83^\circ 11' \cos 61^\circ 30'
    + \sin 83^\circ 11' \sin 61^\circ 30' \cos\angle P \\
  \cos\angle P &= -\cot 83^\circ 11' \cot 61^\circ 30'
    = -0.1195 \times 0.5430 \\
  &= -0.0649 \\
  \angle P &= 93^\circ 43' 16'' = 6^{\mathrm{h}} 14^{\mathrm{m}} 53^{\mathrm{s}} \\
  2\angle P &= 12^{\mathrm{h}} 29^{\mathrm{m}} 46^{\mathrm{s}}
\end{align*}
\hfill(30\%)
